\bigskip
\noindent Autres pistes de l'onglet Findings pour ce thème, dans des cotes qui n'ont pas de section ici :\par

\begin{encadre}
\begin{theoreme*}[Grothendieck, page 12, critère d'immersion ouverte\piste{102-open-immersion-criterion}]\nonverifie
Sur une base localement noethérienne $S$, un monomorphisme formellement étale $u : G \to F$ de foncteurs sur les $S$-schémas est une immersion ouverte dès que $F$ et $G$ sont des faisceaux fpqc qui commutent aux limites inductives filtrantes d'anneaux et aux limites adiques d'anneaux artiniens.
\end{theoreme*}
\end{encadre}
{\small Démonstration donnée par l'édition, non vérifiée en Lean.\par}
{\small La page dit \q{infinitésimalement étale}, et le mot \q{adiques} y est incertain. Mathlib n'a pas la descente fpqc des foncteurs sur les schémas. On se reportera aux notes de Grothendieck (\lot{102}{1}{12}).\par}

\begin{encadre}
\begin{theoreme*}[Grothendieck, pages 2 à 4, 15 et 19, critère valuatif de platitude d'une adhérence\piste{127-closure-flatness-artinian-or-trait}]\nonverifie
Soient $X \to Y$ localement noethérien, $F$ cohérent sur $X$, $U$ un ouvert rétrocompact et $G_U$ un quotient cohérent de $F|_U$, plat sur $Y$. Si les anneaux locaux de $X$ aux points de $X - U$ ont des fibres formelles géométriquement normales, alors $G = \mathrm{Im}(F \to i_*G_U)$ est plat sur $Y$, avec $U$ universellement $G$-dense relativement à $Y$, si et seulement s'il en est ainsi après tout changement de base $Y' \to Y$ où $Y'$ est artinien ou un trait.
\end{theoreme*}
\end{encadre}
{\small Les feuillets en esquissent une démonstration incomplète (pages 3 et 4). La page~4 est en grande partie illisible, et le \q{vrai théorème} de la page~15, qui porte \q{artinien ou trait}, est sans démonstration ; le mot \q{artinien} n'est incertain que dans la remarque~1.4.2 qui le précède. Mathlib n'a pas la densité schématique universelle ni les fibres formelles. On se reportera aux notes de Grothendieck (\lot{127}{1}{2}).\par}

\begin{encadre}
\begin{theoreme*}[Grothendieck, pages 7, 9, 10 et 12, ouverture de Hironaka sur une base réduite ; voir aussi Kollár 2011\piste{127-hironaka-reduced-base-openness}]\nonverifie
Soient $f : X \to Y$ de type fini, $Y$ noethérien, et $\mathcal O_{Y,y}$ réduit à fibres formelles géométriquement normales. Supposons $f$ universellement ouvert en $x$, à fibres équidimensionnelles de même dimension près de $x$, $X_y$ réduit aux générisations maximales de $x$ et $(X_y)_{\mathrm{red}}$ géométriquement normal. Alors $x$ a un voisinage ouvert $U$ tel que $U_{\mathrm{red}} \to Y$ soit plat à fibres géométriquement normales.
\end{theoreme*}
\end{encadre}
{\small Les feuillets en esquissent une démonstration incomplète (pages 9 à 12), avec un \q{lemme à dégager} jamais écrit. D'après le résumé de l'article de Kollár, Hironaka traite le cas d'une base régulière. Mathlib n'a pas les fibres formelles ni les morphismes universellement ouverts. On se reportera aux notes de Grothendieck (\lot{127}{1}{7}).\par}

\begin{encadre}
\begin{theoreme*}[Grothendieck, pages 2 et 3, descente non plate sur une base artinienne\piste{127-non-flat-descent-artinian-epimorphism}]\nonverifie
Sur une base artinienne $Y$, le problème d'un quotient plat $G$ de $F$ qui prolonge $G_U$, avec $G \to i_*G_U$ universellement injectif relativement à $Y$, descend le long de tout épimorphisme de schémas $Y' \to Y$, plat ou non, pourvu que $Y' \to Y$ soit fini, ou que $f$ et la solution soient de présentation finie. La solution, si elle existe, est unique.
\end{theoreme*}
\end{encadre}
{\small Le lemme~1.3 n'est pas démontré, et l'alternative \q{fini ou de présentation finie} vient d'une note au crayon. Mathlib n'a pas la descente le long des épimorphismes de schémas. On se reportera aux notes de Grothendieck (\lot{127}{1}{2}).\par}

\begin{encadre}
\begin{theoreme*}[Grothendieck, pages 1 à 5, 13 et 15, représentabilité par un schéma localement quasi-fini ; voir aussi Murre 1965\piste{131-locally-quasi-finite-representability}]\nonverifie
Un foncteur $F$ sur les schémas au-dessus d'un $S$ localement noethérien est représentable par un $S$-schéma séparé localement quasi-fini si et seulement s'il vérifie la descente fpqc, testée sur $B \to \hat B$, la présentation finie locale, l'effectivité sur $\hat B = \varprojlim B_n$, la pro-représentabilité sur les anneaux artiniens par des anneaux finis sur eux, le critère valuatif de séparation, et deux conditions de \q{modularité}, l'une locale sur les bases locales complètes, l'autre générique.
\end{theoreme*}
\end{encadre}
{\small Les feuillets en esquissent une démonstration incomplète (pages 1 à 5) : la dernière étape s'arrête à la condition de constructibilité, et deux propositions formelles ne sont pas démontrées. Murre traite le cas non ramifié et dit seulement que le cas localement quasi-fini est possible. Mathlib n'a pas les critères de représentabilité de type Artin. On se reportera aux notes de Grothendieck (\lot{131}{1}{1}).\par}

\begin{encadre}
\begin{theoreme*}[Grothendieck, pages 8, 9 et 19, Tor génériques et suites régulières : l'énoncé en un point\piste{94-generic-tor-regular-sequences}]
Soient $A$ un anneau commutatif, $J$ un idéal nilpotent de $A$ ($J^n = 0$), $A_0 = A/J$, et $R$ un $A$-module quelconque. On suppose que les gradués $\mathrm{gr}^k_J(A) = J^kA/J^{k+1}A$ sont des $A_0$-modules libres et que les gradués $\mathrm{gr}^k_J(R) = J^kR/J^{k+1}R$ sont des $A_0$-modules plats. Alors toute suite $A$-régulière $t = (t_1, \ldots, t_r)$ est $R$-régulière, et son image dans $A_0$ est $A_0$-régulière. Si $A$ est local noethérien, il suffit que les $\mathrm{gr}^k_J(A)$ soient plats sur $A_0$ ; si de plus $R$ est de type fini non nul et les $t_i$ dans l'idéal maximal, $t$ est une suite $R$-régulière au sens strict. C'est la situation de $A_{\mathfrak p}$, $\mathfrak p \in U$, dans les lemmes~6 et~16 ; le complexe de Koszul de $t$ résolvant $A/(t)$, on en tire $\mathrm{Tor}_i^A(R, A/(t)) = 0$ pour $i > 0$.
\end{theoreme*}
\end{encadre}
{\small \emph{Démonstration.} Trois outils. (1) Deux sur trois : pour une suite exacte $0 \to M' \to M \to M'' \to 0$ et $t$ faiblement régulière sur $M''$, $t$ est faiblement régulière sur $M'$ si et seulement si elle l'est sur $M$, et $M'/(t)M' \to M/(t)M$ reste injective (c'est l'annulation de $\mathrm{Tor}_1(M'', A/(t))$ que demande la page) ; on le voit élément par élément, par le lemme du serpent. (2) Dévissage : si $t$ est faiblement régulière sur les gradués $\mathrm{gr}^k_J(M)$ pour $k < j$, elle l'est sur $M/J^jM$, par la suite exacte $0 \to \mathrm{gr}^j_J(M) \to M/J^{j+1}M \to M/J^jM \to 0$ et (1) ; avec $J^nM = 0$, elle l'est sur $M$. (3) Platitude : si $\bar t$ est $A_0$-régulière, elle l'est sur tout $A_0$-module plat $N = N \otimes_{A_0} A_0$. Le pas de récurrence du lemme~6 se fait alors ainsi. Supposons $\bar t_1, \ldots, \bar t_k$ déjà $A_0$-régulière, et soit $J^m$ le dernier gradué non nul de $A$ ($J^{m+1} = 0$), libre et non nul sur $A_0$. Par (3) et (2), $(t_1, \ldots, t_k)$ est faiblement régulière sur $A/J^m$, donc, par (1), $J^m/(t)J^m$ s'injecte dans $A/(t)A$, où $t_{k+1}$ est régulier ; il l'est donc sur $J^m/(t)J^m \simeq (A_0/(\bar t))^{(I)}$, $I \neq \emptyset$, et, en restreignant à un facteur, sur $A_0/(\bar t_1, \ldots, \bar t_k)$. Une fois $\bar t$ régulière dans $A_0$, (3) la rend régulière sur chaque $\mathrm{gr}^k_J(R)$ et (2) sur $R$. Comme le dit la note de la page~19, la liberté ne sert, côté module, qu'à travers la platitude, et aucune finitude n'est demandée à $R$. Dans le cas local noethérien, les $\mathrm{gr}^k_J(A)$ sont de type fini sur l'anneau local $A_0$, et plats, donc libres. Lean (\href{\leanblob/Folder94.lean}{\texttt{Folder94.\allowbreak{}lean}}) démontre (1) (\texttt{isWeaklyRegular\_iff\_of\_exact}, \texttt{injective\_mapQ\_of\_exact}), (2) (\texttt{isWeaklyRegular\_quotient\_filt}, \texttt{isWeaklyRegular\_of\_layers}), (3) (\texttt{isWeaklyRegular\_of\_flat}), le pas de récurrence (\texttt{isWeaklyRegular\_quotient\_of\_layers\_free}), l'énoncé (\texttt{isWeaklyRegular\_of\_gr}) et sa forme locale (\texttt{isWeaklyRegular\_of\_gr\_flat\_of\_isLocalRing}, \texttt{isRegular\_of\_gr\_flat\_of\_isLocalRing}).\par}
{\small Ne sont pas démontrés en Lean : le passage de la régularité de $t$ sur $R$ à l'annulation des $\mathrm{Tor}_i$ par le complexe de Koszul (Mathlib n'a ni les complexes de Koszul ni la suite exacte longue des $\mathrm{Tor}$ ; seul l'analogue de $\mathrm{Tor}_1$ ci-dessus l'est), et l'ouvert dense de l'énoncé global ci-dessous. On se reportera aux notes de Grothendieck (\lot{94}{1}{8}).\par}

\begin{encadre}
\begin{theoreme*}[Grothendieck, pages 8 et 9, l'ouvert dense pour $X = S$\piste{94-generic-tor-regular-sequences}]
Soient $A$ un anneau noethérien, $J$ un idéal nilpotent de $A$ tel que $A_0 = A/J$ soit réduit (par exemple le nilradical), et $R$ un $A$-module de type fini. Il existe un ouvert dense $U \subset \operatorname{Spec} A$ tel que, pour tout $\mathfrak p \in U$, toute suite $A_{\mathfrak p}$-faiblement régulière d'éléments de $A_{\mathfrak p}$ soit $R_{\mathfrak p}$-faiblement régulière. On peut prendre pour $U$ le lieu, dans $\operatorname{Spec} A = \operatorname{Spec} A_0$, où les gradués non nuls, en nombre fini, $\mathrm{gr}^k_J(A)$ et $\mathrm{gr}^k_J(R)$ sont libres sur $A_0$.
\end{theoreme*}
\end{encadre}
{\small \emph{Démonstration.} Les gradués sont des $A_0$-modules de type fini, donc de présentation finie, et leur lieu de liberté est ouvert ; l'intersection $U$ de ces lieux contient les idéaux premiers minimaux, où l'anneau réduit $A_0$ se localise en un corps, et un ensemble qui contient les premiers minimaux est dense. Les gradués commutent à la localisation : $\mathrm{gr}^k_J(S^{-1}M) = S^{-1}\mathrm{gr}^k_J(M)$, car $J^k(S^{-1}M)$ est le localisé de $J^kM$ ; si $\mathfrak p$ est au-dessus de $\mathfrak q \in \operatorname{Spec} A_0$, on a donc $\mathrm{gr}^k_J(M_{\mathfrak p}) = (A_0)_{\mathfrak q} \otimes_{A_0} \mathrm{gr}^k_J(M)$, libre sur $(A_0)_{\mathfrak q}$ pour $\mathfrak p \in U$. Le pas de récurrence de l'énoncé en un point vaut avec n'importe quelle $A$-algèbre $B$ au lieu de $A_0$ et n'importe quel module non nul $T$ au lieu de $A$, par la même démonstration ; avec $B = (A_0)_{\mathfrak q}$ et $T = A_{\mathfrak p}$, toute suite d'éléments de $A$ faiblement régulière sur $A_{\mathfrak p}$ l'est sur $R_{\mathfrak p}$. Enfin tout élément de $A_{\mathfrak p}$ est le produit d'une unité par l'image d'un élément de $A$, et la régularité faible ne dépend de la suite qu'à des unités près. Lean (\href{\leanblob/Folder94.lean}{\texttt{Folder94.\allowbreak{}lean}}) démontre la commutation à la localisation (\texttt{isLocalizedModule\_layerMap}, \texttt{layerEquiv}), le pas de récurrence généralisé (\texttt{isWeaklyRegular\_of\_gr'}) et l'énoncé (\texttt{exists\_isOpen\_dense\_isWeaklyRegular'}).\par}

\begin{encadre}
\begin{theoreme*}[Grothendieck, pages 8, 9 et 19, Tor génériques et suites régulières : l'ouvert dense\piste{94-generic-tor-regular-sequences}]\nonverifie
Soient $S = \operatorname{Spec} A$ noethérien, non nécessairement réduit, et $R$ un faisceau cohérent sur un $S$-schéma $X$ de type fini. Il existe un ouvert dense $U \subset S$ tel que $\mathrm{Tor}_i^A(R, A_{\mathfrak p}/(t)) = 0$ pour tout $i > 0$, tout $\mathfrak p \in U$ et toute suite $A_{\mathfrak p}$-régulière~$t$. C'est un substitut de la platitude générique, qui tombe en défaut sur une base non réduite.
\end{theoreme*}
\end{encadre}
{\small Le pas de récurrence du lemme~6, que la page dit \q{bien standard}, est de l'édition, et la page écrit $\mathrm{gr}_J(\mathcal O_X)$ là où l'argument demande $\mathrm{gr}_J(R)$. L'énoncé en un point et le cas $X = S$ ci-dessus, pour la régularité des suites, sont vérifiés en Lean ; restent le passage aux $\mathrm{Tor}$ par le complexe de Koszul, et, pour $X$ de type fini, la platitude générique, que Mathlib n'a pas. On se reportera aux notes de Grothendieck (\lot{94}{1}{8}).\par}

\begin{encadre}
\begin{theoreme*}[Grothendieck, pages 15 et 16, critère de représentabilité sur $\mathrm{Mod}(A)$ ; voir aussi Hall 2013\piste{94-representability-criterion}]\nonverifie
Soit $A$ un anneau noethérien quotient d'un anneau régulier. Un foncteur $F : \mathrm{Mod}(A) \to \mathrm{Ens}$ est de la forme $\mathrm{Hom}_A(P, -)$ avec $P$ fini si et seulement s'il est exact à gauche, commute aux limites inductives filtrantes, vérifie une condition de complétion $\mathfrak m_{\mathfrak p}$-adique et une $g$-adique sur les modules finis, envoie les modules finis sur des modules finis, a une fonction $\mathfrak p \mapsto \dim F(k(\mathfrak p))$ constructible et, pour tout quotient $B$ et tout $B$-module fini $\Omega$, commute génériquement à la réduction de $\Omega$ modulo les systèmes de paramètres. La conservation des produits n'en fait pas partie.
\end{theoreme*}
\end{encadre}
{\small Les feuillets en esquissent une démonstration incomplète (pages 15 et 16) : plusieurs liens sont de l'édition, et la démonstration telle qu'elle est écrite n'est sûre que pour $A$ excellent. Hall énonce un autre critère, par la conservation des produits. Mathlib n'a pas ces conditions de complétion sur les foncteurs de modules. On se reportera aux notes de Grothendieck (\lot{94}{1}{15}).\par}

\noindent Piste dont l'énoncé dépend d'une lecture à vérifier :
\begin{itemize}\small\setlength{\itemsep}{0pt}
\item Cote~94 : pour $A$ noethérien dont tout quotient intègre est génériquement normal, et $F$ exact à gauche, commutant aux limites inductives filtrantes et vérifiant la condition de complétion locale et la finitude locale, la représentabilité de $F$ par un module fini sur un ouvert non vide de chaque $V(\mathfrak p)$ équivaut à la constructibilité de $\mathfrak p \mapsto \dim_{k(\mathfrak p)} F(k(\mathfrak p))$. La réciproque repose sur le milieu de la page~7, en grande partie illisible, et l'hypothèse de normalité sur une lecture incertaine, à vérifier sur le fac-similé. Mathlib n'a pas la constructibilité de ces fonctions de rang. On se reportera aux notes de Grothendieck (\lot{94}{1}{7}).\piste{94-rank-constructibility}
\end{itemize}

\begin{encadre}
\begin{contrexemple}[Page 7, la réciproque de la remarque 5 bis sous la finitude locale]\verifie[\leanblobmain]{Folder94}
Sur $A = k[x, y]$, $k$ un corps infini, soient $a_2, a_3, \dots$ des éléments distincts de $k$ et $h_j = x - a_j$. Le foncteur $F(M) = \varinjlim\,(M \xrightarrow{h_2} M \xrightarrow{h_3} M \to \cdots)$ est l'identité sur chaque $A_{\mathfrak p}$, vérifie la condition de complétion et la finitude locale, et $\dim_{k(\mathfrak p)} F(k(\mathfrak p)) = 1$ pour tout $\mathfrak p$ ; pourtant $F(A_g)$ n'est de type fini sur $A_g$ pour aucun $g \neq 0$, et $F$ n'est représentable sur aucun ouvert non vide de $\operatorname{Spec} A$.
\end{contrexemple}
\end{encadre}
{\small \emph{Démonstration.} Un idéal premier contient au plus un $h_j$, donc presque tous les $h_j$ sont inversibles dans $A_{\mathfrak p}$ et $F$ y est l'identité à partir d'un rang (\texttt{tele\_atPrime\_bijective}) ; d'où les conditions locales et la dimension~$1$ des fibres (\texttt{finrank\_tele\_field}). Un $g \neq 0$ n'est divisible que par un nombre fini de $h_j$, et $F(A_g)$ est la réunion strictement croissante des $(h_2 \cdots h_i)^{-1} A_g$ (\texttt{not\_finite\_tele\_away}, \texttt{not\_representable\_away}). L'exactitude à gauche et la commutation aux limites inductives filtrantes sont classiques et ne sont pas formalisées. Sous la finitude globale, $F(M)$ de type fini pour tout $M$ de type fini, que l'esquisse de la page utilise, la réciproque tient, et sans l'hypothèse sur le lieu normal : sur un ouvert de $V(\mathfrak p)$ où le rang vaut $n$, les modules locaux de la proposition~4 sont libres de rang $n$, $F(B)$ est projectif de type fini et $F(B) \otimes M \to F(M)$ est un isomorphisme en chaque point. Les deux étapes d'algèbre commutative de cet argument sont démontrées en Lean ci-dessous ; la proposition~4 et le passage du foncteur $F$ à ces énoncés ne le sont pas. Quelle finitude la page écrit au juste reste à lire sur le fac-similé.\par}

\begin{encadre}
\begin{theoreme*}[Page 7, rang constant sur un anneau réduit\piste{94-rank-constructibility}]
Soient $B$ un anneau local réduit, de corps résiduel $k$, et $P$ un $B$-module de type fini tel que $\dim_{B_{\mathfrak q}} (B_{\mathfrak q} \otimes_B P) = \dim_k (k \otimes_B P)$ pour tout idéal premier minimal $\mathfrak q$ ($B_{\mathfrak q}$ est alors un corps). Alors $P$ est libre. Par suite, si $P$ est de présentation finie sur un anneau $B$, à localisés libres aux idéaux maximaux, et si $\varphi \colon P \otimes_B M \to N$ est bijectif après localisation en chaque idéal maximal, $P$ est projectif et $N \simeq \operatorname{Hom}_B(P^\vee, M)$.
\end{theoreme*}
\end{encadre}
{\small \emph{Démonstration.} Une base de $k \otimes P$ se relève, par Nakayama, en une surjection $B^n \to P$. En un premier minimal $\mathfrak q$, elle devient une surjection entre deux $B_{\mathfrak q}$-espaces de même dimension $n$, donc une bijection ; les coordonnées d'un élément du noyau ont donc une image nulle dans chaque $B_{\mathfrak q}$, sont dans chaque premier minimal, donc nilpotentes, donc nulles. Aucune hypothèse noethérienne n'est utile. Pour le second point, la projectivité se lit sur les localisés en les idéaux maximaux, de même que la bijectivité de $\varphi$, et $P \otimes M \simeq P^{\vee\vee} \otimes M \simeq \operatorname{Hom}(P^\vee, M)$ pour $P$ projectif de type fini. C'est ce qu'on applique à $P = F(B)$ et $\varphi \colon F(B) \otimes M \to F(M)$ dans la réciproque sous la finitude globale. Lean (\href{\leanblob/Folder94.lean}{\texttt{Folder94.\allowbreak{}lean}}) : \texttt{free\_of\_finrank\_eq\_of\_isReduced}, \texttt{nonempty\_equiv\_hom\_dual}.\par}
